\section{习题}

\subsection{基础过关题}

\begin{enumerate}

\item \textbf{（单项选择题）} 主机 H 通过本地域名服务器 L 解析域名 \texttt{www.example.com}，L 的缓存为空。下列叙述中正确的是
\begin{center}
\begin{tabular}{>{\raggedright\arraybackslash}p{0.45\linewidth}>{\raggedright\arraybackslash}p{0.45\linewidth}}
(A) H 对 L 采用递归查询时，L 必须返回最终结果或失败，不能只给出下一台服务器地址 &
(B) H 对 L 采用迭代查询时，L 必须代 H 完成后续全部查询 \\[6pt]
(C) 迭代查询中，根服务器会代替 L 继续向 \texttt{.com} 服务器查询 &
(D) 递归查询中，根服务器若没有该记录，可以只返回 \texttt{.com} 服务器地址
\end{tabular}
\end{center}

\ifshowsolutions\par\noindent\textbf{解析：} 递归查询要求被查询的服务器返回最终结果或失败信息，不能只返回“去问某台服务器”的转介，故 (A) 正确；但 L 自己向其他服务器查询时仍可采用迭代方式，这两件事并不矛盾。(B) 把递归的义务错误地安到了迭代上：迭代查询只负责“指路”，给出下一台服务器的地址，由请求者自己继续查询。(C) 中“代查到底”正是递归查询的行为，不是迭代。(D) 与递归的定义矛盾，递归服务器不能只给转介。\boxed{\textbf{答案：}(A)}\else \answerarea \fi

\item \textbf{（单项选择题）} 下列关于 DNS 层次结构和区的说法中，正确的是
\begin{center}
\begin{tabular}{>{\raggedright\arraybackslash}p{0.45\linewidth}>{\raggedright\arraybackslash}p{0.45\linewidth}}
(A) 一个域只能划分为一个区 &
(B) 本地域名服务器属于 DNS 层次结构的一级 \\[6pt]
(C) 区可以小于或等于它所属的域 &
(D) 根域名服务器保存所有二级域名的 IP 地址
\end{tabular}
\end{center}

\ifshowsolutions\par\noindent\textbf{解析：} 域是名字空间中的一棵子树，区是某台权限域名服务器实际管辖的范围。一个域可以划分成多个区（把子域委派给下级服务器管理），而一个区只属于一个域，因此区小于或等于域，(C) 正确，(A) 错误。本地域名服务器是主机访问 DNS 的代理，不属于“根—顶级—权限”的层次结构，(B) 错误。根服务器只保存各顶级域的委派信息（以及少量根区数据），并不保存全部二级域名的 IP 地址，(D) 错误。\boxed{\textbf{答案：}(C)}\else \answerarea \fi

\item \textbf{（单项选择题）} 关于 FTP 的两条 TCP 连接，下列说法中正确的是
\begin{center}
\begin{tabular}{>{\raggedright\arraybackslash}p{0.45\linewidth}>{\raggedright\arraybackslash}p{0.45\linewidth}}
(A) 控制连接由服务器主动发起 &
(B) 主动模式下，服务器从端口 20 向客户端告知的端口发起数据连接 \\[6pt]
(C) 被动模式下，服务器从端口 21 向客户端发起数据连接 &
(D) 数据连接在整个会话期间保持，控制连接每传完一个文件就关闭
\end{tabular}
\end{center}

\ifshowsolutions\par\noindent\textbf{解析：} FTP 使用两条 TCP 连接。控制连接由客户端发起并连到服务器的 21 端口，在整个会话期间保持，用于传送 FTP 命令和应答；数据连接每传输一个文件或一份目录列表时单独建立，传完即关闭，故 (A)、(D) 错误。主动模式（PORT）下，客户端告知自己的监听端口，服务器从自己的 20 端口主动连接该端口，(B) 正确。被动模式（PASV）下，服务器打开一个临时端口并告知客户端，由客户端发起数据连接，服务器不会从 21 端口发起，(C) 错误。\boxed{\textbf{答案：}(B)}\else \answerarea \fi

\item \textbf{（单项选择题）} 用户 A 把一封邮件发送给用户 B：(1) A 的用户代理把邮件提交给 A 的邮件服务器；(2) A 的邮件服务器把邮件转发给 B 的邮件服务器；(3) B 用 IMAP 客户端从自己的邮件服务器读取邮件。第 (1)(2)(3) 步所用协议依次是
\begin{center}
\begin{tabular}{ll}
(A) SMTP、SMTP、IMAP & (B) SMTP、IMAP、SMTP \\[6pt]
(C) IMAP、SMTP、SMTP & (D) SMTP、SMTP、SMTP
\end{tabular}
\end{center}

\ifshowsolutions\par\noindent\textbf{解析：} SMTP 负责“发送”和邮件服务器之间的中转：第一步，用户代理向自己的邮件服务器提交邮件用 SMTP；第二步，两台邮件服务器之间转发邮件也用 SMTP；第三步，用户代理从接收方邮件服务器读取邮件属于“接收”，应使用 IMAP 或 POP3，题干指定 IMAP 客户端，故第三步为 IMAP。SMTP 不能用于读取邮箱中的邮件。\boxed{\textbf{答案：}(A)}\else \answerarea \fi

\item \textbf{（单项选择题）} 关于 DHCP 的工作过程，下列说法中正确的是
\begin{center}
\begin{tabular}{>{\raggedright\arraybackslash}p{0.45\linewidth}>{\raggedright\arraybackslash}p{0.45\linewidth}}
(A) 客户端以 0.0.0.0 为源、255.255.255.255 为目的广播 DHCPDISCOVER &
(B) DHCPREQUEST 总是单播发送，以便其他服务器保留已预留的地址 \\[6pt]
(C) 所有收到 DHCPDISCOVER 的服务器都必须发送 DHCPOFFER &
(D) 分配到的地址永久有效，续租时必须重新执行完整四步
\end{tabular}
\end{center}

\ifshowsolutions\par\noindent\textbf{解析：} DHCP 四步为 DHCPDISCOVER $\to$ DHCPOFFER $\to$ DHCPREQUEST $\to$ DHCPACK。客户端此时还没有可用 IP，DISCOVER 用源地址 0.0.0.0、目的地址 255.255.255.255（受限广播），承载在 UDP 上，客户端用 68 端口发送、服务器在 67 端口接收，故 (A) 正确。REQUEST 采用广播，是为了让未被选中的服务器得知自己的 OFFER 未被接受并释放预留，故 (B) 错误。服务器可以不应答 DISCOVER，OFFER 也可单播或广播，故 (C) 错误。地址有租用期，租期过半时客户端可直接向原服务器发送 REQUEST 续租，无需重走完整四步，故 (D) 错误。\boxed{\textbf{答案：}(A)}\else \answerarea \fi

\item \textbf{（填空题）} 在 HTTP/1.1 中，(1) 请求首部 \underline{\qquad\qquad} 用于指定目标主机，使一台服务器可以区分多个虚拟主机；(2) 客户端缓存副本仍然有效时，服务器返回状态码 \underline{\qquad\qquad}；(3) HTTP/1.1 默认使用 \underline{\qquad\qquad} 连接。

\ifshowsolutions\par\noindent\textbf{解析：} HTTP/1.1 要求请求报文必须带 \texttt{Host} 首部，服务器据此选择虚拟主机；条件请求中若缓存副本仍然有效，服务器返回 304 Not Modified，响应不携带实体主体；HTTP/1.1 默认使用持久连接，同一 TCP 连接可依次传送多个对象，只有显式发送 \texttt{Connection: close} 才在本次响应后关闭。\boxed{\text{答案：}\text{(1) Host}\quad\text{(2) 304}\quad\text{(3) 持久}}\else \answerarea \fi

\end{enumerate}

\subsection{真题风格与综合训练}

\begin{enumerate}

\item \textbf{（单项选择题）} 某网页包含 1 个基本 HTML 文件和 5 个图像，全部位于同一服务器。忽略 DNS 时间和对象传输时间，且必须先收到基本 HTML 才能知道 5 个图像的 URL。非持久串行、持久非流水线、持久流水线三种方式所需的 RTT 数依次为
\begin{center}
\begin{tabular}{ll}
(A) $12$、$7$、$3$ & (B) $12$、$8$、$3$ \\[6pt]
(C) $10$、$7$、$3$ & (D) $12$、$7$、$2$
\end{tabular}
\end{center}

\ifshowsolutions\par\noindent\textbf{解析：} 共 1 个 HTML 加 5 个图像，即 6 个对象。非持久连接下，每个对象都要单独建立 TCP 连接并用 1 RTT 完成请求—响应，串行共 $6\times 2=12$ RTT。持久非流水线只需建立一次连接（1 RTT），取 HTML 用 1 RTT，再逐个取 5 个图像用 5 RTT，共 $1+1+5=7$ RTT。持久流水线在收到 HTML 后一次发出全部 5 个请求，响应连续返回，图像部分只占 1 RTT，共 $1+1+1=3$ RTT。\boxed{\textbf{答案：}(A)}\else \answerarea \fi

\item \textbf{（单项选择题）} 关于 FTP 与电子邮件协议，下列说法中正确的是
\begin{center}
\begin{tabular}{>{\raggedright\arraybackslash}p{0.45\linewidth}>{\raggedright\arraybackslash}p{0.45\linewidth}}
(A) FTP 的控制连接和数据连接都贯穿整个会话 &
(B) FTP 数据连接可传二进制文件；SMTP 传二进制附件需 MIME 编码 \\[6pt]
(C) POP3 提供服务器端多文件夹和多设备同步，IMAP 不提供 &
(D) SMTP 的 MAIL FROM 与邮件首部 From 都表示邮件作者
\end{tabular}
\end{center}

\ifshowsolutions\par\noindent\textbf{解析：} FTP 只有控制连接贯穿会话，数据连接每传一个文件或一份目录列表就建立一次、传完即关闭，故 (A) 错误。FTP 的数据连接是透明的字节流，可直接传二进制文件；SMTP 原始协议按 7 位 ASCII 设计，二进制附件必须经 MIME 规定的 Base64 或 Quoted-Printable 等传输编码，故 (B) 正确。POP3 以把邮件下载到本地为主，不提供 IMAP 式的服务器端多文件夹和多设备状态同步，(C) 把两者说反了。SMTP 的 \texttt{MAIL FROM} 是信封上的发件人，用于投递和退信；邮件首部 \texttt{From} 是显示给收件人的作者，两者可以不同，故 (D) 错误。\boxed{\textbf{答案：}(B)}\else \answerarea \fi

\item \textbf{（真题风格，综合题）} 某网页由 1 个基本 HTML 文件（$20$ KB）和 4 幅 JPEG 图像（每幅 $100$ KB）组成，全部放在同一台服务器上。设往返时间 $\mathrm{RTT}=100$ ms，链路带宽为 $10$ Mb/s，为便于计算取 $1$ KB $=1000$ B；忽略 DNS 解析时间、报文首部开销和连接释放时间，对象的传输时间按“对象大小/带宽”在 RTT 之外另行计算。必须先收到基本 HTML 才能知道 4 幅图像的 URL；建立 TCP 连接用 1 RTT，连接建立后一次请求—响应用 1 RTT。分别计算：
(1) 非持久连接、串行方式的总时间；
(2) 持久连接、非流水线方式的总时间；
(3) 持久连接、流水线方式的总时间；
(4) 若客户端在建立 TCP 连接前就已获得 4 幅图像的 URL，可在连接建立后立即把 5 个请求一起流水线发出，总时间为多少？与 (3) 相差多少？四种方式中哪种最短？

\ifshowsolutions\par\noindent\textbf{解析：}
\textbf{思路：} 每种方式的时间都由三部分组成：新建 TCP 连接数 $\times 1$ RTT、请求—响应次数 $\times 1$ RTT、所有对象的传输时间之和。先把传输时间算出来：
\[
T_{\text{HTML}}=\frac{20\times1000\times8}{10\times10^{6}}\ \mathrm{s}=16\ \mathrm{ms},
\qquad
T_{\text{img}}=\frac{100\times1000\times8}{10\times10^{6}}\ \mathrm{s}=80\ \mathrm{ms}.
\]
(1) 非持久连接下 5 个对象各建一次连接、各做一次请求—响应：
\[
T_1=5\times(2\times100)+16+4\times80=1000+16+320=1336\ \mathrm{ms}.
\]
(2) 持久非流水线只建立一次连接，HTML 与 4 幅图像依次请求：
\[
T_2=100+(100+16)+4\times(100+80)=100+116+720=936\ \mathrm{ms}.
\]
(3) 流水线方式下，收到 HTML 后一次性发出 4 个图像请求，图像阶段只等待 1 个 RTT，传输时间仍为 $4\times80$：
\[
T_3=100+(100+16)+(100+4\times80)=100+116+420=636\ \mathrm{ms}.
\]
(4) 建连 1 RTT 后即可连续发出 5 个请求，不必等 HTML 到达；首个响应仍需 1 RTT，全部数据的传输时间为 $16+4\times80=336$ ms，故
\[
T_4=100+100+336=536\ \mathrm{ms}.
\]
与 (3) 相比少 $100$ ms，正好是“必须先收到 HTML 才能知道图像 URL”所多付的 1 RTT；四种方式中 (4) 最短。
\boxed{\textbf{答案：}(1)\,1336\ \mathrm{ms}\quad(2)\,936\ \mathrm{ms}\quad(3)\,636\ \mathrm{ms}\quad(4)\,536\ \mathrm{ms}}\else \answerarea[10] \fi

\item \textbf{（真题风格，综合题）} 某主机新接入一个局域网，尚未配置 IP 地址。该局域网内有 DHCP 服务器和本地域名服务器，出口路由器是该局域网的默认网关；本地域名服务器无缓存，对上级服务器采用迭代查询。主机随后访问 \texttt{http://www.example.com/index.html}，该 Web 服务器位于另一个网络。回答下列问题：
(1) 写出 DHCP 获取地址的四个报文名称及顺序；指出其中由客户端广播发送的报文；写出 DHCPDISCOVER 的源 IP 地址、目的 IP 地址，以及服务器和客户端使用的 UDP 端口。
(2) 写出本地域名服务器完成解析的迭代查询过程（向谁查询、得到什么），它共发送几条查询报文？主机向本地域名服务器发送几条查询报文？若在 TTL 内另一台主机查询同一域名，本地域名服务器如何处理？
(3) 主机向 Web 服务器发送 IP 数据报前，需要用 ARP 获得哪个设备的 MAC 地址？为什么不是 Web 服务器的 MAC 地址？该 ARP 请求的作用范围如何？
(4) 从主机发出 TCP SYN 到收到 HTTP 响应，最少需要多少 RTT（不计 DNS 时间和传输时间）？HTTP/1.1 默认采用哪种连接方式？

\ifshowsolutions\par\noindent\textbf{解析：}
(1) 四步为 DHCPDISCOVER $\to$ DHCPOFFER $\to$ DHCPREQUEST $\to$ DHCPACK。客户端的 DISCOVER 和 REQUEST 采用广播：DISCOVER 时主机还没有 IP，源地址为 $0.0.0.0$，目的地址为受限广播地址 $255.255.255.255$，客户端从 UDP 68 端口发往服务器的 67 端口；REQUEST 广播是为了让其他提供过 OFFER 的服务器释放预留。OFFER 和 ACK 可以单播，也可以广播。

(2) 本地域名服务器依次向三台服务器查询：第一步，向根域名服务器查询，得到 \texttt{.com} 顶级域名服务器的地址；第二步，向 \texttt{.com} 顶级域名服务器查询，得到 \texttt{example.com} 权限域名服务器的地址；第三步，向该权限域名服务器查询，得到 \texttt{www.example.com} 的 A 记录（IP 地址）。三次查询都只拿到“下一站”或最终记录，由本地域名服务器自己继续，因此共发送 3 条查询报文。主机只向本地域名服务器发送 1 条递归查询报文。若在 TTL 内另一台主机查询同一域名，本地域名服务器直接用缓存中的记录回答，不再向上级发送查询。

(3) 需要获得默认网关（出口路由器本网接口）的 MAC 地址。目的 IP 位于另一个网络，主机不能用 ARP 直接解析跨网段的 Web 服务器地址；它只能把 IP 数据报封装成帧发给默认网关，由网关继续转发，所以帧的目的 MAC 必须是网关的 MAC。ARP 请求在本局域网内广播，网关收到后以单播方式返回自己的 MAC 地址。

(4) 从发出 SYN 起算：1 RTT 完成三次握手，再用 1 RTT 完成 HTTP 请求—响应，共 2 RTT。HTTP/1.1 默认使用持久连接，同一 TCP 连接可继续传送后续对象。
\boxed{\text{答案：}\text{(1) D-O-R-A}\quad\text{(2) 3 条}\quad\text{(3) 网关 MAC}\quad\text{(4) 2 RTT}}\else \answerarea[10] \fi

\end{enumerate}
